Throughout, $T = 1/8\pi M$ and $S = 4\pi M^{2} = A/4$ are those of Theorem \claimref{synthesis.quarter-from-temperature}, so $TS = M/2$ and $\beta = 8\pi M$.

\textbf{(i).} For a Gibbs state $\varrho = e^{-\beta H}/Z$ with $e^{-\beta H}$ of trace class, $-\ln\varrho = \beta H + \ln Z$, so $S(\varrho) = -\tr\varrho\ln\varrho = \beta\langle H\rangle + \ln Z$ (on finite blocks this is the Gibbs form of Theorem \claimref{singularity.kms-uniqueness-factor}). With $\langle H_M\rangle = M$ and $S = A/4$: $\ln Z_M = A/4 - \beta M = 4\pi M^{2} - 8\pi M^{2} = -4\pi M^{2} = -A/4$, and $F := -T\ln Z_M = M/2 = M - TS$. This is the Gibbons-Hawking value: their Euclidean action of the Schwarzschild section is $I = \beta M/2 = 4\pi M^{2}$ and $\ln Z = -I$ (G.\,W.~Gibbons and S.\,W.~Hawking, Phys.\ Rev.\ D 15 (1977) 2752, cited from memory); the content here is only its rewriting in the series' variables.

\textbf{(ii).} Since $H_M = E_0 + H'_M$, $Z_M = e^{-\beta E_0}Z'_M$ and $\langle H_M\rangle = E_0 + \langle H'_M\rangle$. The Gibbs identity for $H'_M$ (the same state) gives $S = \beta\langle H'_M\rangle + \ln Z'_M$. As $H'_M \geq 0$ with $0$ in its spectrum and $e^{-\beta H'_M}$ trace class, $0$ is an eigenvalue of finite multiplicity $g_0 \geq 1$ and $Z'_M = g_0 + \sum_{\lambda > 0}g_\lambda e^{-\beta\lambda} \geq 1$, with equality only if the spectrum is the single nondegenerate eigenvalue $0$; then $\varrho_M$ is pure and $S = 0$, contradicting $S = A/4 > 0$. Hence $Z'_M > 1$, $\ln Z'_M > 0$, and $\beta\langle H'_M\rangle = S - \ln Z'_M < S$, that is $\langle H'_M\rangle < TS = M/2$; so $E_0 = M - \langle H'_M\rangle > M/2$, while $\langle H'_M\rangle \geq 0$ gives $E_0 \leq M$. Equivalently $E_0 = F + T\ln Z'_M > F = M/2$. The same identity applied to any sector whose Hamiltonian has ground energy $0$, in a Gibbs state at $\beta$, gives energy $< T\cdot(\text{its entropy})$; with entropy at most $A/4$ this is $< M/2$. For one fermion mode of weight $\omega$ it is the inequality $\beta\omega n < s(x)$ of Theorem \claimref{synthesis.residue-from-mode-entropy}(v), $s(x) - x\,n(x) = \ln(1 + e^{-x}) > 0$ by Theorem \claimref{synthesis.boundary-mode-thermal-entropy}. In the symmetric ordering $\omega(N - \tfrac12)$ the mode carries $\beta E = x(n - \tfrac12)$; with $u = e^{-x}$, $n - \tfrac12 = -(1 - u)/(2(1 + u)) = -\tfrac12\tanh\tfrac{x}{2}$, so $\beta E = -\tfrac{x}{2}\tanh\tfrac{x}{2} \leq 0 < s(x)$; the change of ordering is the constant shift $E_0 \mapsto E_0 - \omega/2$, which leaves $\langle H_M\rangle$, $S$ and the bounds on $E_0$ for the given $H_M$ unchanged.

\textbf{(iii).} $E_0 = M$ if and only if $\langle H'_M\rangle = 0$; since $\varrho_M$ is faithful on the range of $e^{-\beta H_M}$ (every eigenvector has positive weight), $\tr(\varrho_M H'_M) = 0$ with $H'_M \geq 0$ forces $H'_M = 0$, so $H_M = E_0$ on the support. For $H_M = E_0 + \sum_j\omega_jv_j^{*}v_j$ on the Fock space of $J$, the Gibbs state is the product of the single-mode Gibbs states (the marginals of Theorem \claimref{synthesis.boundary-mode-thermal-entropy}), $\langle H'_M\rangle = E = \sum_j\omega_jn_j$ and $E_0 = M - E$. By Theorem \claimref{synthesis.residue-from-mode-entropy}(ii) the equality case $\kappa_c = 1/(4\ln 2)$ is all $\omega_j = 0$, which is $H'_M = 0$, $E = 0$, $E_0 = M$, $Z'_M = 2^{|J|}$, $S = |J|\ln 2$; then $\ln Z_M = -\beta E_0 + |J|\ln 2 = -S$ gives $\beta E_0 = 2|J|\ln 2 = 2S$, the Smarr relation $M = 2TS$ for the ground energy. At positive weights $Z'_M = \prod_j(1 + e^{-x_j}) > 1$ and $E > 0$, so $M/2 < E_0 < M$.

\textbf{(iv).} By (ii), $M/2 < E_0(M) \leq M$ for every $M$, so $E_0(M) \to 0$ as $M \to 0$ and $E_0(M) \to \infty$ as $M \to \infty$; no constant satisfies both bounds on $(0, \infty)$. For a fixed $H$ with $e^{-\beta H}$ trace class on an interval of $\beta$, $\frac{d}{d\beta}\langle H\rangle_\beta = -(\langle H^{2}\rangle_\beta - \langle H\rangle_\beta^{2}) \leq 0$ (differentiate $\tr(He^{-\beta H})/\tr e^{-\beta H}$), whereas along the family $M = \beta/8\pi$ increases with $\beta$; so no fixed Hamiltonian has $\langle H\rangle_\beta = M(\beta)$ on any interval, and $H_M$ depends on $M$. This is the standard negative heat capacity of the Schwarzschild family (Hawking, 1976, cited from memory), $dM/dT = -1/(8\pi T^{2})$, read as a statement about the Hamiltonian: the dependence on $M$ cannot be absorbed into a fixed dynamics and is a dependence on the state's label.

\textbf{Check.} The script verifies, in exact rational-function arithmetic in $u = e^{-\beta\varepsilon}$, that $Z' > 1$ on $(0, 1)$ for spectra with more than one state (one fermion mode, the Fock space of weights $1, 2$, a truncated oscillator, a degenerate ground) and $Z' = 1$ for a single nondegenerate level; the variance identity $u\,d\langle H\rangle/du = \operatorname{Var}(H)$ with $\operatorname{Var} > 0$ on $(0, 1)$, which is $d\langle H\rangle/d\beta = -\varepsilon\operatorname{Var} < 0$; the zero-weight sector $Z' = 2^{|J|}$, $\langle H'\rangle = 0$ for $|J| = 1, \ldots, 6$ with $e^{-\beta E_0} = (e^{-S})^{2} = 4^{-|J|}$; and $Z' = (1 + u)(1 + u^{2})$, $\langle H'\rangle > 0$ at weights $1, 2$; and the symmetric-ordering identity $n - \tfrac12 = -(1 - u)/(2(1 + u))$ with $n - \tfrac12 < 0$ on $(0, 1)$. The analytic step (the Gibbs identity) is in the text; positivity on $(0, 1)$ is certified by Bernstein coefficients.
